% Options for packages loaded elsewhere
\PassOptionsToPackage{unicode}{hyperref}
\PassOptionsToPackage{hyphens}{url}
\documentclass[
  11pt,
]{article}
\usepackage{xcolor}
\usepackage[left=3cm,right=3cm,top=2.5cm,bottom=2.5cm]{geometry}
\usepackage{amsmath,amssymb}
\setcounter{secnumdepth}{5}
\usepackage{iftex}
\ifPDFTeX
  \usepackage[T1]{fontenc}
  \usepackage[utf8]{inputenc}
  \usepackage{textcomp} % provide euro and other symbols
\else % if luatex or xetex
  \usepackage{unicode-math} % this also loads fontspec
  \defaultfontfeatures{Scale=MatchLowercase}
  \defaultfontfeatures[\rmfamily]{Ligatures=TeX,Scale=1}
\fi
\usepackage{lmodern}
\ifPDFTeX\else
  % xetex/luatex font selection
\fi
% Use upquote if available, for straight quotes in verbatim environments
\IfFileExists{upquote.sty}{\usepackage{upquote}}{}
\IfFileExists{microtype.sty}{% use microtype if available
  \usepackage[]{microtype}
  \UseMicrotypeSet[protrusion]{basicmath} % disable protrusion for tt fonts
}{}
\usepackage{setspace}
\makeatletter
\@ifundefined{KOMAClassName}{% if non-KOMA class
  \IfFileExists{parskip.sty}{%
    \usepackage{parskip}
  }{% else
    \setlength{\parindent}{0pt}
    \setlength{\parskip}{6pt plus 2pt minus 1pt}}
}{% if KOMA class
  \KOMAoptions{parskip=half}}
\makeatother
\usepackage{longtable,booktabs,array}
\usepackage{caption}
\captionsetup[table]{skip=6pt}
\newcounter{none} % for unnumbered tables
\usepackage{calc} % for calculating minipage widths
% Correct order of tables after \paragraph or \subparagraph
\usepackage{etoolbox}
\makeatletter
\patchcmd\longtable{\par}{\if@noskipsec\mbox{}\fi\par}{}{}
\makeatother
% Allow footnotes in longtable head/foot
\IfFileExists{footnotehyper.sty}{\usepackage{footnotehyper}}{\usepackage{footnote}}
\makesavenoteenv{longtable}
\usepackage{graphicx}
\makeatletter
\newsavebox\pandoc@box
\newcommand*\pandocbounded[1]{% scales image to fit in text height/width
  \sbox\pandoc@box{#1}%
  \Gscale@div\@tempa{\textheight}{\dimexpr\ht\pandoc@box+\dp\pandoc@box\relax}%
  \Gscale@div\@tempb{\linewidth}{\wd\pandoc@box}%
  \ifdim\@tempb\p@<\@tempa\p@\let\@tempa\@tempb\fi% select the smaller of both
  \ifdim\@tempa\p@<\p@\scalebox{\@tempa}{\usebox\pandoc@box}%
  \else\usebox{\pandoc@box}%
  \fi%
}
% Set default figure placement to htbp
\def\fps@figure{htbp}
\makeatother
\setlength{\emergencystretch}{3em} % prevent overfull lines
\providecommand{\tightlist}{%
  \setlength{\itemsep}{0pt}\setlength{\parskip}{0pt}}
%% tools/stamp.tex
%% Provenance footer for PDFs rendered from this project.
%% Vendored by zzcollab; this file is not project-specific. The
%% footer prints the source document and the time of rendering.
%% The git version is defined here too, and so is available to a
%% project that wants it on the page, but it is left off the
%% default footer: it already names the staged copy in share/ and
%% appears in its manifest row. All values are supplied at render
%% time by tools/stamp-render.R, which writes a small generated
%% file that redefines the macros below. The \providecommand
%% defaults keep a bare LaTeX compile working when the document is
%% built without the wrapper.
\usepackage{fancyhdr}
\usepackage{xcolor}
\providecommand{\stampsource}{\detokenize{(source not recorded)}}
\providecommand{\stamptime}{(time not recorded)}
\providecommand{\stampversion}{\detokenize{(version not recorded)}}
\setlength{\footskip}{40pt}
\newcommand{\stampfootcontent}{%
  \footnotesize\color{gray}%
  \parbox[b]{\textwidth}{%
    \stamptime\hfill\thepage\\[2pt]%
    \ttfamily\raggedright\stampsource}}
\pagestyle{fancy}
\fancyhf{}
\fancyfoot[C]{\stampfootcontent}
\renewcommand{\headrulewidth}{0pt}
\renewcommand{\footrulewidth}{0.4pt}
%% The title page uses the `plain` page style; redefine it so the
%% stamp appears there too.
\fancypagestyle{plain}{%
  \fancyhf{}%
  \fancyfoot[C]{\stampfootcontent}%
  \renewcommand{\headrulewidth}{0pt}%
  \renewcommand{\footrulewidth}{0.4pt}}
\renewcommand{\stampsource}{\textasciitilde{}/\allowbreak{}prj/\allowbreak{}res/\allowbreak{}36-pmsimstats-ng/\allowbreak{}pmsimstats-ng/\allowbreak{}analysis/\allowbreak{}report/\allowbreak{}01-dgp-mean-moderation-vs-mvn/\allowbreak{}bullets.Rmd}
\renewcommand{\stamptime}{2026-08-15 16:47 PDT}
\renewcommand{\stampversion}{\detokenize{bc0b671-wip}}
\usepackage{amsmath}
\usepackage{amssymb}
\usepackage{booktabs}
\usepackage{longtable}
\usepackage{float}
\usepackage[mathlines]{lineno}
\linenumbers
\usepackage{bookmark}
\IfFileExists{xurl.sty}{\usepackage{xurl}}{} % add URL line breaks if available
\urlstyle{same}
% fallback for those not using the hyperref driver hyperxmp:
\makeatletter
\@ifundefined{xmpquote}{\newcommand{\xmpquote}[1]{#1}}{}
\makeatother
\hypersetup{
  pdftitle={Structured summary: mean moderation and covariance moderation as data-generating architectures for biomarker-treatment interactions in aggregated N-of-1 trials (supplementary key-points summary)},
  pdfauthor={pmsimstats team},
  hidelinks,
  pdfcreator={LaTeX via pandoc}}

\title{Structured summary: mean moderation and covariance moderation as data-generating architectures for biomarker-treatment interactions in aggregated N-of-1 trials (supplementary key-points summary)}
\author{pmsimstats team}
\date{August 15, 2026}

\begin{document}
\maketitle

{
\setcounter{tocdepth}{2}
\tableofcontents
}
\setstretch{1.1}
\emph{This document is a structured, section-by-section summary of
the key points argued in the companion manuscript, ``Mean moderation
and covariance moderation as data-generating architectures for
biomarker-treatment interactions in aggregated N-of-1 trials, and
their divergent power under carryover.'' It is provided as a
supplementary reading aid and does not stand on its own; all
notation, derivations, simulation detail, and citations are in the
main manuscript.}

\bigskip\noindent

\rule{\textwidth}{0.4pt}\bigskip

\section{1. Introduction}

\textbf{Power for these designs is estimated by simulation, and the
simulation must specify the biomarker interaction in its DGP in
more than one possible way.}

\begin{itemize}\setlength{\itemsep}{2pt}\setlength{\parskip}{0pt}
\item The target of inference is the biomarker-treatment
interaction.
\item Closed-form power exists only in simplified cases
[@pmsimstats-paper08]. For the designs and model considered here it
is obtained by Monte Carlo simulation.
\item How the interaction is built into the DGP is a modeling
choice, and that choice is the subject of this paper.
\end{itemize}

\textbf{There are two architectures, mean moderation and
differential correlation.}

\begin{itemize}\setlength{\itemsep}{2pt}\setlength{\parskip}{0pt}
\item Under mean moderation an interaction term scales the
treatment effect in the conditional mean. This is standard across
trial-simulation frameworks.
\item Under covariance moderation the biomarker changes the
response-outcome correlation, so the interaction is represented in
the joint distribution. This is the approach taken by Hendrickson et
al. for N-of-1 trials.
\item A dual-channel architecture activating both at once is
developed in a companion paper (Section 2.2).
\end{itemize}

\textbf{Hendrickson et al. chose covariance moderation for the
prazosin and PTSD trial as the biologically appropriate choice,
departing from the modal mean-moderation convention. This paper
quantifies what that choice costs under carryover.}

\begin{itemize}\setlength{\itemsep}{2pt}\setlength{\parskip}{0pt}
\item Covariance moderation was the logical specification given the
trial's own results, because mean moderation would have imposed a
deterministic biomarker-response relationship the biology did not
support.
\item Mean moderation remains the modal specification in the N-of-1
and broader biomarker-simulation literature, so this was a
deliberate departure, not an oversight.
\item The size of the resulting power penalty under carryover had
not previously been quantified. Doing so is this paper's central
aim.
\end{itemize}

\textbf{Aggregated N-of-1 trials pool repeated within-patient
on-drug and off-drug contrasts, and carryover is the nuisance that
distinguishes them.}

\begin{itemize}\setlength{\itemsep}{2pt}\setlength{\parskip}{0pt}
\item A single N-of-1 trial repeatedly crosses one patient between
treatment states, and pooling many such series yields
population-level inference at small $N$.
\item Carryover is residual drug effect that decays over days to
weeks and contaminates off-drug occasions.
\item Three within-subject designs, namely crossover (CO),
open-label with blinded discontinuation (OL+BDC), and Hybrid,
differ in how on-drug and off-drug occasions are arranged.
\end{itemize}

\section{2. Two DGP architectures}

\textbf{Notation. $B_i$ is the biomarker, $D_{it}$ the on-drug
indicator, $BR_{it}$ the biological response, and \texttt{Dbc} the
analysis-layer exposure.}

\begin{itemize}\setlength{\itemsep}{2pt}\setlength{\parskip}{0pt}
\item $B_i$, $D_{it}$, $BR_{it}$ describe the data-generating layer.
\item The analysis layer uses a continuous exposure \texttt{Dbc}
(Section 3.1).
\item Typewriter identifiers denote code and formula variables.
Mathematical symbols are used otherwise.
\end{itemize}

\subsection{2.1 Mean moderation}

\textbf{Under mean moderation the interaction is explicit in the
conditional mean, so that a higher biomarker value gives a
proportionally larger treatment effect.}

\begin{itemize}\setlength{\itemsep}{2pt}\setlength{\parskip}{0pt}
\item $\beta_{bm}$ is the biomarker moderation parameter on the
centered biomarker.
\item The effect scales with $B_i$ deterministically in the
conditional mean.
\item The signal lives in the first moment of the outcome.
\end{itemize}

\textbf{The multiplicative form maps onto the standard additive
effect-modeling regression, so both are instances of mean
moderation.}

\begin{itemize}\setlength{\itemsep}{2pt}\setlength{\parskip}{0pt}
\item The additive interaction $\beta_3$ equals
$\beta_1 \cdot \beta_{bm}$ in the multiplicative form.
\item The additive form adds a prognostic main effect $\beta_2$ the
multiplicative form sets to zero.
\item For interaction power the two are interchangeable up to that
prognostic distinction.
\end{itemize}

\subsection{2.2 Covariance moderation through differential correlation}

\textbf{Under covariance moderation the interaction is in the
second moment. The biomarker modulates BR through a correlation
that decays during the off-drug period.}

\begin{itemize}\setlength{\itemsep}{2pt}\setlength{\parskip}{0pt}
\item The conditional mean of BR rises with $b$ through the
correlation, not through a mean parameter.
\item The effective correlation is $c_{bm}$ on drug and $c_{bm}
e^{-\lambda t_{sd}}$ off drug.
\item Moderation persists with attenuated strength off drug and
vanishes only at complete decay.
\end{itemize}

\section{3. Consequences for Statistical Power Under Carryover}

\subsection{3.1 Empirical demonstration}

\textbf{Under mean moderation the carryover-induced power loss is
modest in the two-period designs (2-3\% relative) and larger in the
Hybrid design (7\%).}

\begin{itemize}\setlength{\itemsep}{2pt}\setlength{\parskip}{0pt}
\item Relative loss to $t_{1/2} = 1$: 1.9\% (CO, $z = 0.71$),
6.9\% (Hybrid, $z = 3.34$), 2.8\% (OL+BDC, $z = 1.07$).
\item Only the Hybrid loss is distinguishable from zero. CO and
OL+BDC are within Monte Carlo error.
\item At or below the 8-13\% range earlier exploratory runs
suggested.
\end{itemize}

\textbf{Under covariance moderation the loss is much larger and strongly
design-dependent, reaching 30.6\% at OL+BDC but absent at CO.}

\begin{itemize}\setlength{\itemsep}{2pt}\setlength{\parskip}{0pt}
\item Relative loss: 3.0\% (CO, not significant), 15.4\% (Hybrid),
30.6\% (OL+BDC), roughly 11 times the matched
mean-moderation loss.
\item The CO cell fluctuates within MC error, consistent with no
architecture-by-carryover interaction there.
\item Below the 40-60\% range earlier runs suggested, but confirming
substantial, design-dependent erosion.
\end{itemize}

\subsection{3.2 Why the architectures diverge}

\textbf{Under mean moderation, carryover only compresses the
treatment contrast, and the biomarker signal-to-noise ratio is
nearly preserved.}

\begin{itemize}\setlength{\itemsep}{2pt}\setlength{\parskip}{0pt}
\item Carryover shrinks \texttt{Dbc} and inflates off-drug BR means
(shared with covariance moderation).
\item But $\beta_{bm}$ is a fixed mean parameter with drug-state-
independent noise.
\item So the interaction variance falls but the signal-to-noise
ratio degrades only modestly.
\end{itemize}

\textbf{Covariance moderation's interaction signal is attacked from both
sides, an attack that mean moderation escapes, which accounts for the
divergence.}

\begin{itemize}\setlength{\itemsep}{2pt}\setlength{\parskip}{0pt}
\item The \texttt{bm:Dbc} coefficient depends on the covariance of
$B_i$ with \texttt{Dbc}-weighted outcomes.
\item Under covariance moderation both the \texttt{Dbc} contrast and
the BM-BR correlation decay together.
\item Under mean moderation the second attack is absent, so power is
largely preserved.
\end{itemize}

\subsection{3.3 Robustness check at $N = 140$}

\textbf{An $N = 140$ confirmation run checks that the architecture
ordering survives larger samples.}

\begin{itemize}\setlength{\itemsep}{2pt}\setlength{\parskip}{0pt}
\item Section 3.1 puts $N = 70$ on this design, where the loss is
mechanically visible (no ceiling saturation).
\item The check runs 12 cells at $N = 140$ (70 per path), 800
replicates each, on the two-path OL+BDC design.
\item 100\% convergence across 9{,}600 fits.
\end{itemize}

\textbf{At $N = 140$ the ordering between the two architectures
persists, though absolute magnitudes are compressed by ceiling
saturation.}

\begin{itemize}\setlength{\itemsep}{2pt}\setlength{\parskip}{0pt}
\item Covariance moderation loses 9.3 points ($z = 7.4$), and mean
moderation loses 1.7 points ($z = 2.0$).
\item Both sit within 2-5 points of the ceiling at $t_{1/2} = 0$,
truncating the absorbable loss.
\item The Section~3.1 results remain the canonical reference for
the narrative magnitude.
\end{itemize}

\section{4. Which recipe matches which biology?}

\textbf{Use mean moderation when the biomarker mechanically sets the
size of the drug effect.}

\begin{itemize}\setlength{\itemsep}{2pt}\setlength{\parskip}{0pt}
\item The biomarker fixes effective dose or receptor occupancy, so
it scales the effect proportionally.
\item Doubling the biomarker roughly doubles the effect at every
timepoint.
\item Carryover keeps the same proportion in the off-drug period, so power is
barely touched.
\end{itemize}

\textbf{Use covariance moderation when the biomarker is only a proxy for a
hidden responder state.}

\begin{itemize}\setlength{\itemsep}{2pt}\setlength{\parskip}{0pt}
\item The biomarker imperfectly indexes the real cause, so identical
biomarkers can respond differently.
\item It predicts response only while the drug is acting.
\item In the off-drug period the correlation fades and power drops sharply, as
in Section 3.
\end{itemize}

\textbf{The motivating prazosin/PTSD case is a covariance-moderation
story, so its carryover power loss is a clinically relevant
warning.}

\begin{itemize}\setlength{\itemsep}{2pt}\setlength{\parskip}{0pt}
\item Biomarker is resting blood pressure, a proxy for noradrenergic
tone, not a driver of drug levels.
\item Designs with long off-drug periods may lose more power than a
mean-moderation simulation predicts.
\item If blood pressure also acts through the mean channel, the
dual-channel architecture of the companion paper applies.
\end{itemize}

\section{5. Implications for Simulation Study Design}

\subsection{5.1 Choosing the appropriate architecture}

\textbf{When the mechanism is ambiguous, use the
covariance moderation as the conservative default.}

\begin{itemize}\setlength{\itemsep}{2pt}\setlength{\parskip}{0pt}
\item Covariance moderation is the safer default for design decisions.
\item Its larger carryover sensitivity keeps power estimates from
being optimistic.
\item A dual-channel sweep across mixing weights is available in
the companion paper.
\end{itemize}

\subsection{5.2 Reporting requirements}

\textbf{These requirements specialize the ADEMP template, and
the architecture choice is its data-generating (D) element.}

\begin{itemize}\setlength{\itemsep}{2pt}\setlength{\parskip}{0pt}
\item ADEMP fixes Aims, Data-generating mechanisms, Estimands,
Methods, and Performance measures.
\item The DGP-architecture choice is the D element.
\item Our reporting rules are ADEMP specialized to
architecture-sensitive biomarker simulations.
\end{itemize}

\section{6. Can a smarter analysis recover the lost power?}

\textbf{Covariance moderation's loss has two sources, and only
the variance-compression part is recoverable, by down-weighting
low-signal off-drug points.}

\begin{itemize}\setlength{\itemsep}{2pt}\setlength{\parskip}{0pt}
\item Source 1: the drug exposure variable's variance shrinks
(recoverable).
\item Source 2: the biomarker-response correlation genuinely fades
off-drug (lost information).
\item Methods that down-weight or drop off-drug points can help with
the first.
\end{itemize}

\section{7. How this fits the wider literature}

\textbf{Mean moderation is near-universal in the literature, so
published power estimates may be optimistic for covariance moderation
biomarkers.}

\begin{itemize}\setlength{\itemsep}{2pt}\setlength{\parskip}{0pt}
\item Simon, Freidlin-Korn, Renfro-Sargent, and standard N-of-1
studies all put the effect in the mean.
\item Even N-of-1 carryover work sharing Hendrickson's authors uses
mean moderation, so Hendrickson's choice looks like a one-off.
\item Most published power figures may be optimistic for genuinely
covariance moderation biomarkers.
\end{itemize}

\textbf{The distinction matters most in treatment-switching
designs. When the mechanism is unknown, compute power both ways
and use the lower number.}

\begin{itemize}\setlength{\itemsep}{2pt}\setlength{\parskip}{0pt}
\item Treatment-switching designs (crossover, N-of-1, basket) put
patients through off-drug phases where the covariance-moderation correlation decays.
\item Parallel-group designs have no off-drug phase, so the
architecture choice does not matter.
\item This is a DGP question, separate from the analysis-model
guidance of the PATH statement.
\end{itemize}

\section{References}\label{references}

\protect\phantomsection\label{refs}

\end{document}
